\begin{abstract}
This paper addresses a single, frozen question: which enterprise states
are logically decidable, and which are provably undecidable, under
incomplete observability when enterprises are modeled as continua
$K_{EA}$ within the Ontology of Continua (OC Core) and its executable
specification ETS.K\_EA.v1.1.

Decidability is treated strictly as a logical property of state
predicates, not as an algorithmic, statistical, inferential, or
procedural notion. Observability is formalized exclusively through the
ETS-defined observation operator and observable signature space, and is
explicitly separated from decidability. No inference, reconstruction,
probabilistic assumptions, or control semantics are introduced.

We show that an enterprise state predicate is decidable if and only if
it is invariant over the observational equivalence classes induced by
admissible observations. Predicates that depend on internal structural
features not fixed by observable signatures are shown to be undecidable
by construction. As a central negative result, we prove that most
non-terminal phase distinctions cannot be decided under incomplete
observability.

A sharp boundary case is identified in the form of \emph{STOP-only
decidability}: under explicit OC/ETS conditions, STOP can be certified
as a terminal outcome even when all non-STOP phase distinctions remain
undecidable. This yields a complete trichotomy of enterprise state
predicates into decidable, undecidable, and STOP-only decidable classes.

All results are fully internal to the OC/ETS formalism and are
explicitly falsifiable at the level of observable signatures. The paper
makes no claims about estimation, governance, optimization, or
intervention, and establishes formal limits on what can and cannot be
decided about enterprise states from admissible observations alone.
\end{abstract}

\noindent\textbf{Keywords:} enterprise continuum; partial observability;
logical decidability; indistinguishability; diagnosability; STOP.
